\chapter{Implementation and Results}\label{sec:ImplementationAndResults}\label{Chapter4} % For referencing the chapter elsewhere, use \ref{Chapter4}
% Briefly restate the research objectives and the purpose of this chapter.
% Provide an overview of the implementation strategy and key findings.
% Outline how this chapter is structured.
This chapter presents the implementation details and key findings of the research system designed
to enhance static malware analysis using \glspl{llm} and \gls{rag}. It outlines the architectural
realization of the proposed pipelines, describes how individual components were integrated, and
evaluates their performance across a representative corpus of malicious \glspl{pe} samples. The
experiments were conducted on a dataset of 15 malware samples across five commonly studied
families~\cite{berrios2025malware,abrazak2016rise,mat2021systematic} (See
section~~\ref{subsec:MalwareSamplesSourcing}), enabling comparative evaluation of classification
behavior.

\section{Environment}\label{sec:Environment}
% Describes the Python-based toolchain used across the project:
% - **PDM** for managing dependencies and isolating virtual environments.
% - **Malwarebazaar/VX-Underground/Malpedia**: Data sources for the RAG datastore
% - **PostgreSQL + pgvector** as the vector store for document embeddings.
% - **Haystack** for RAG orchestration.
% - **Ghidra + PyGhidra** for static feature extraction from PE files.
% - **OpenAI API** for LLM inference.
% - **VirusTotal API** for metadata enrichment.
% - Versioning and reproducibility strategies.
The experimental environment was built with open-source tools and APIs, with an emphasis on modular
components, reproducibility, and transparent workflows. All development was carried out in
\gls{python}, using containerized services and reproducible dependency management.

\vspace{0.5em}
\noindent This overview describes the environment and supporting tools used across
the project. A step-by-step guide to the setup and configuration can be found in the
repository:~\footnote{\url{https://github.com/andremmfaria/rexis}}

\vspace{1em}
\noindent The main tools and frameworks are:
\begin{itemize}
      \item \textbf{\gls{python}} --~Primary implementation language for orchestration logic and
            pipeline components. Python was chosen for its strong ecosystem in cybersecurity,
            data analysis, and \gls{machine-learning}.\footnote{\url{https://www.python.org}}

      \item \textbf{PDM (Python Development Master)} --~Manages dependencies and virtual
            environments. It provides deterministic builds and isolated execution, supporting
            reproducibility across runs and machines.\footnote{\url{https://pdm-project.org}}

      \item \textbf{\gls{postgresql} + \gls{pgvector}} --~Relational database extended with the
            \texttt{pgvector} module\footnote{\url{https://www.postgresql.org},
                  \url{https://github.com/pgvector/pgvector}} for similarity search on embeddings. It
            acts as the vector store backing the retrieval stage of the \gls{rag} pipeline. The
            database runs in a Docker container for consistency and portability.\footnote{\url{https://www.docker.com}}

      \item \textbf{\gls{haystack}} --~Framework used to implement the \gls{rag} workflow,
            including ingestion, chunking, embedding, and retrieval. Its modular components
            simplify pipeline orchestration and integration with external
            \glspl{llm}.\footnote{\url{https://haystack.deepset.ai}}

      \item \textbf{\gls{gpt4o} (inference-only)} --~No fine-tuning was
            performed. The model is used strictly for contextual reasoning and static
            feature analysis, relying on structured prompts enriched with retrieved
            context.\footnote{\url{https://platform.openai.com/docs/api-reference/introduction}}

      \item \textbf{\gls{ghidra} + PyGhidra} --~\gls{ghidra} serves as the static analysis and
            decompilation tool. PyGhidra enables automated extraction of imports, functions,
            strings, and decompiled code from \gls{pe} samples.\footnote{\url{https://github.com/NationalSecurityAgency/ghidra},
                  \url{https://github.com/NationalSecurityAgency/ghidra/blob/master/Ghidra/Features/PyGhidra/src/main/py}}

      \item \textbf{Docker} --~Used to containerize services such as \gls{postgresql} with
            \gls{pgvector}. Containerization ensures reproducibility, isolation, and consistent
            builds across environments.\footnote{\url{https://www.docker.com}}

      \item \textbf{\gls{malwarebazaar} API} --~Source of real-world samples, hashes, tags, and
            metadata, used for enrichment and validation.\footnote{\url{https://bazaar.abuse.ch/api/}}

      \item \textbf{\gls{virustotal} API} --~Supplies scan metadata, antivirus tags, and vendor
            consensus labels. This data supports the baseline classifier and reconciliation
            logic.\footnote{\url{https://developers.virustotal.com/reference}}
\end{itemize}
Version control was handled with Git and GitHub, ensuring traceability of code and
configuration changes. All pipelines were developed to run deterministically under PDM allowing
results to be reproduced on different systems with the same setup.

\section{Implementation Overview}\label{sec:ImplementationOverview}
% Describe the technical implementation of the system and pipelines.
% Cover how the REXIS CLI is structured (e.g., modules for ingestion, classification).
% Include information on:
%   - Folder structure and entry points (e.g., Typer CLI commands, src/rexis layout)
%   - Configuration management (e.g., Dynaconf or environment variables)
%   - Integration of tools (Ghidra, PostgreSQL, Haystack, LLM APIs)
%   - How components interact during runtime (feature extraction → retrieval → prompt → inference)
%   - Use of Docker or containerization for reproducibility
% Optionally include short code snippets, CLI command samples, or log outputs.
The system is implemented as a modular Python application using the
\texttt{Typer}\footnote{\url{https://typer.tiangolo.com/}} framework for CLI orchestration. The
codebase is organized to support two primary analysis pipelines: a baseline static classifier and
an enhanced \gls{llm}-driven pipeline backed by \gls{rag}. Common modules handle preprocessing,
data ingestion, feature extraction, and logging, while each pipeline defines its own classification
logic and output formatting.

\subsection{Common Modules}\label{subsec:CommonModules}
The following components are shared across all pipeline operations:

\begin{itemize}
      \item \textbf{Sample Loader:} Loads \gls{pe} files from the designated
            corpus directory and calculates hashes (SHA256, MD5) for sample tracking.

      \item \textbf{Ghidra Integration:} Static feature extraction is performed
            using \gls{ghidra} via the \texttt{PyGhidra} API. A containerized Ghidra server
            is initialized at runtime, and samples are submitted in batch mode.

      \item \textbf{Feature Extractor:} Extracts structural and semantic features,
            including PE headers, DLL imports, printable strings, section entropy, and
            entry point offsets. Results are serialized into both JSON and \texttt{pandas}
            DataFrames for consistency and reuse.

      \item \textbf{Logger:} Captures structured logs for each pipeline execution,
            including runtime events, errors, hashes, timestamps, and tool versions.
            Logging is centralized for reproducibility and debugging.

      \item \textbf{Configuration Manager:} Runtime parameters, including API
            keys, credentials, and feature toggles, are managed via
            \texttt{dynaconf}\footnote{\url{https://www.dynaconf.com/}} to allow
            profile-specific overrides and secure separation of secrets.

      \item \textbf{CLI Interface:} A unified CLI under the \texttt{rexis} command,
            implemented using the \texttt{Typer}\footnote{\url{https://typer.tiangolo.com/}} framework,
            enables isolated or batch execution of both pipelines via subcommands
            (e.g., \texttt{rexis run baseline}, \texttt{rexis run rag}).
\end{itemize}

\subsection{Data Ingestion Pipeline}\label{subsec:DataIngestionImpl}
The ingestion pipeline populates the \gls{rag} vector database by processing external threat
intelligence sources into semantically searchable embeddings.

\begin{itemize}
      \item \textbf{Document Loader:} Gathers raw documents from predefined
            directories, API clients (e.g., \gls{malwarebazaar}), and web archives
            (e.g., \gls{malpedia}, VX-Underground). Each document is tagged with
            source metadata and a unique identifier.

      \item \textbf{Preprocessing Engine:} Cleans input text by removing
            boilerplate, embedded scripts, and redundant entries. It performs
            language normalization and token cleaning.

      \item \textbf{Chunker:} Splits documents into semantically meaningful
            sections using length-aware heuristics (e.g., newline or heading-based
            segmentation). This step ensures each chunk remains within the token
            window limits of the embedding model.

      \item \textbf{Embedder:} Encodes each chunk using a transformer-based
            embedding model, such as OpenAI's \texttt{text-embedding-3-small}.
            Embeddings are batched and streamed to reduce latency.

      \item \textbf{Vector Store Writer:} Inserts chunks and associated metadata
            into a \gls{postgresql} + \gls{pgvector} database. The Haystack
            \texttt{DocumentStore} abstraction is used to manage indexing,
            similarity metrics, and metadata-based filtering.
\end{itemize}

The ingestion process is automated via the CLI using the command \texttt{rexis ingest
      --source=<name>} and may be run incrementally. All ingested data is versioned to ensure
reproducibility and traceability of semantic retrieval during classification.

\subsection{Baseline Pipeline Implementation}\label{subsec:BaselineImplementation}
The baseline pipeline uses deterministic rule-based logic combined with metadata enrichment via
\gls{virustotal}. Its architecture includes the following modules:

\begin{itemize}
      \item \textbf{Heuristic Classifier:} Applies manually curated rules to
            extracted features. For instance, the presence of obfuscation-related
            imports or unusually high section entropy can trigger detection of
            suspicious behavior.

      \item \textbf{VirusTotal Enricher:} Queries the \gls{virustotal} API using
            the sample's hash and retrieves engine detections, tags, and metadata.
            Results are normalized into canonical family labels via a tag-mapping layer.

      \item \textbf{Decision Engine:} Merges heuristic and VirusTotal outputs
            using a confidence-based reconciliation strategy. If conflicts occur,
            priority is determined by a configurable threshold rule.

      \item \textbf{Output Formatter:} Produces structured classification
            summaries, which include heuristic scores, final labels, and audit logs.
\end{itemize}

\subsection{LLM+RAG Pipeline Implementation}\label{subsec:RAGImplementation}
The enhanced pipeline performs semantic retrieval and inference through external \glspl{llm},
enabling context-aware classification and explanation.

\begin{itemize}
      \item \textbf{Semantic Query Builder:} Transforms extracted features into
            natural language questions (e.g., 'What malware family uses these API
            calls and strings?'), incorporating string samples, header attributes, and
            import function names.

      \item \textbf{Retriever:} Uses Haystack's retriever module to search the
            \gls{pgvector} database for top-$k$ relevant documents. Each retrieval includes
            document ID, similarity score, and context snippet.

      \item \textbf{Prompt Builder:} Generates a structured prompt containing the
            static features, retrieved content, and task-specific instruction
            templates. This step uses a fixed template with content placeholders.

      \item \textbf{Compliance Filter:} Scans and sanitizes the final prompt to
            remove potentially harmful or non-compliant artifacts (e.g., large hex
            dumps, file paths, sensitive domains), mitigating prompt injection risks.

      \item \textbf{Inference Client:} Sends the prompt to either OpenAI’s
            \gls{gpt4o} or \gls{deepseek-r1} via their respective APIs. Responses
            are logged with latency, token usage, and engine identifiers.

      \item \textbf{Result Parser:} Extracts structured information from the
            LLM response, including malware classification, rationale, and references to
            retrieved documents.
\end{itemize}

\section{Summary of Key Findings}\label{sec:SummaryOfFindings}
% Summarize the main results obtained from both the baseline and the LLM+RAG pipelines.
% Avoid detailed numerical analysis — focus on broad outcomes such as:
%   - LLM+RAG's classification accuracy compared to the baseline
%   - Quality of explanations or reasoning given by the LLM
%   - Retrieval relevance and diversity across malware families
%   - Support for flexible input types (e.g., script-based malware)
% Highlight patterns and trends that emerged across experiments.
TODO...

\section{Interpretation of Findings}\label{sec:InterpretationOfFindings}
% Analyze and interpret the significance of the results in relation to your research questions.
% Discuss:
%   - How RAG improved static analysis (accuracy, relevance, interpretability)
%   - Where the LLM pipeline underperformed and possible causes
%   - Degree of alignment between LLM inferences and VirusTotal tags
%   - Evidence of hallucinations or compliance issues
% Cross-reference previous literature (e.g., Fujii 2024, Gao 2023) to contextualize your outcomes.
% Discuss whether findings were expected, and any surprises or anomalies encountered.
TODO...

\section{Implications of the Study}\label{sec:Implications}
% Reflect on the practical and academic significance of your research.
% Discuss how your system design and findings may:
%   - Influence future tools for static malware analysis
%   - Support security teams through explainable classification
%   - Extend to other cybersecurity domains (e.g., phishing, vulnerability triage)
% Consider both short-term and long-term implications for real-world use.
TODO...

\section{Limitations}\label{sec:Limitations}
% Clearly state the limitations of your system and experiments.
% Address issues such as:
%   - Limited sample diversity or dataset size
%   - Scope limited to static analysis, excluding behavioral/dynamic data
%   - Heuristic rules being brittle or oversimplified
%   - Dependence on third-party APIs (e.g., OpenAI, VirusTotal)
%   - No adversarial robustness or user-study evaluation
% Be transparent and balanced — link limitations to specific design choices where applicable.
TODO...

\section{Recommendations for Future Research}\label{sec:FutureWork}
% Based on your experience and findings, propose:
%   - Future exploration of hybrid pipelines (static + dynamic)
%   - Inclusion of larger, labeled datasets
%   - Alternative LLMs or fine-tuned models for specific tasks
%   - Improvements in retrieval quality (e.g., better chunking, embeddings)
%   - Real-world evaluation or deployment (e.g., SOC integration, red-team use)
% Ensure recommendations are feasible and rooted in what you observed during implementation.
TODO...

\section{Conclusion}\label{sec:ChapterConclusion}
% Recap the chapter by:
%   - Summarizing what was implemented and what results were achieved
%   - Highlighting how your work addresses the research objectives
%   - Providing a closing reflection on the value of combining LLMs and RAG for static malware analysis
%   - Transitioning into the final conclusion chapter
TODO...
